\section{Day 36: Schr\"odinger in Higher Dimensions (Mar. 9, 2026)}
We skipped 9.3 in Strauss, but you should read it. We looked at \( u_{t} = k \Delta u \) for \( k > 0 \). We found the fundamental solution \( x \in \mathbb{R}^{d}\), \(t > 0 \) whose Kernel is
\[ S_{t}(x) = \frac{1}{\left( 4 \pi k t \right)^{d/2}} e^{- \frac{|x|^{2}}{4 k t}}. \]
Then we moved on to Schr\"odinger very briefly. We will expound upon this. Specifically, the plan is
\begin{enumerate}

	\item Use analogy to the formula with the heat.
	\item Do some comparison to the heat equation. There are many differences.
	\item Solve the harmonic oscillator.

\end{enumerate}
The Schr\"odinger equation with one electron is
\[ - i \hbar u_{t} = \frac{\hbar}{2 m} \Delta u + \frac{\ell^{2}}{r} u, \]
where \( \hbar \) is Planck's constant, and \( \ell \) is a constant, and \( r \) is the radius. We're going to change coordinates here to get the equation
\[ -i u_{t} = \frac{1}{2} \Delta u + V(x) u. \]
Here, \( V \) represents the potential somehow.

Let \( u \colon I \times U \to \mathbb{C} \). We say \( u = \operatorname{Re}(u) + i \operatorname{Im}(u) \). Then if we do \( u_{t} = k \Delta u \) for \( k \in \mathbb{C} \) with \( \operatorname{Re} k > 0 \), this it the proof of the heat equation we had before. So instead, let's look at
\[ u_{t} = \frac{i + \varepsilon}{2} \Delta u \quad \varepsilon > 0. \]
This is not entirely rigorous, but look at this formula and take the limit as \( \varepsilon \to 0 \) to get our fundamental solution. The solution to
\[ \left\{\begin{NiceMatrix}u_{t} = \frac{i}{2} \Delta u = 0 \\ u(0, x) = f(x)\end{NiceMatrix}\right.  \]
is given by
\[ u(t, x) = \frac{1}{\left( 2 \pi i t \right)^{d/2}} \int_{\mathbb{R}^{d}} e^{- \frac{|x-y|^{2}}{2it}} f(y) \, dy. \]
The meaning of \( \sqrt{i} \) here is just \( \lim_{\varepsilon \to 0} \sqrt{i + \varepsilon} = \frac{1+i}{\sqrt{2}} \). 
\begin{remark}
	This is very different from the heat equation. This is mainly due to the \( i \) in the denominator of our exponent of \( e \). Let's see some ways in which it's very different for fun.
\end{remark}

\begin{proposition}
	\begin{itemize}
	
		\item The Schr\"odinger equation is time reversible. If \( u(t, x) \) is a solution, then \( \bar u(-t, x) \) is a solution.
		\item Conservation of total mass. So long as the integral is defined,
			\[ \int_{\mathbb{R}^{d}} |u(t, x)|^{2} \, dx \equiv c. \]
			What's the physical meaning of Schr\"odinger? Well \( |u(t, x)|^{2} \) represents the probability density to find our particle at position \( x \) at time \( t \). I don't know this might be wrong, I'm not a physicist. Naturally then this mass should be conserved for this to make sense, otherwise the probability of the particle existing will change which is ridiculous.
		\item We also see that energy and the momentum are conserved.
			\[ \int_{\mathbb{R}^{d}}  | \nabla u (t, x)|^{2} \, dx, \quad \operatorname{Im} \left( \int_{\mathbb{R}^{d}} \nabla u(t, x) \cdot \bar u(t, x) \, dx \right). \]
	\end{itemize}
\end{proposition}
in the homework or exam, he might ask ``look at \( u_{t} \) plus more terms, is the energy or momentum conserved, what happens?''

\begin{proof}
	Let's start with the mass. If \( M \) is our mass, then
	\begin{align*}
		\frac{d}{dt} M &= \int_{\mathbb{R}^{d}} \partial_{t}u \, \bar u + u \partial_{t} \, \bar u \, dx = 2 \int_{\mathbb{R}^{d}}  \operatorname{Re}(\partial_{t} u \, \bar u) \, dx \\ 
									 &= 2 \operatorname{Re} \left(  \int_{\mathbb{R}^{d}} \frac{i}{2} \Delta u \, \bar u \, dx \right) \\
									 &= 2 \operatorname{Re} \int_{\mathbb{R}^{d}} \int_{\mathbb{R}^{d}}  i\,  \nabla u \,\overline{\nabla u} \, dx \\
									 &= 0  
	\end{align*}

	Now let's check the momentum.
	\begin{align*}
		\partial_{t} \int_{\mathbb{R}^{d}}  \partial_{x_{1}} u \, \bar u \, dx &= \int_{\mathbb{R}^{d}} \partial_{x_{1}} \partial_{t} u \, \bar u	+ \partial_{x_{1}} u \, \overline{\partial_{t} u} \, dx \\
		&= \frac{1}{2} \int_{\mathbb{R}^{d}} \partial_{x_{1}} \, i\Delta \, \bar u - \partial_{x_{1}} u \, i\Delta \bar u \, dx\\
		&= \frac{1}{2} \int_{\mathbb{R}^{d}} - \partial_{x_{1}} \, i \nabla u \, \nabla \bar u \, dx- \frac{1}{2}  \int_{\mathbb{R}^{d}} \partial_{x_{1}} u \, i\Delta \bar u\, dx \\
		&= \frac{1}{2} \int_{\mathbb{R}^{d}} i \partial_{x_{1}} u \, \Delta \bar u \, dx - \frac{1}{2}  \int_{\mathbb{R}^{d}} \partial_{x_{1}} u \, i \Delta \bar u \, dx \\
		&= 0.
	\end{align*}
\end{proof}
As mentioned above, we can look at what happens with \( u_{t} = \frac{i}{2} \Delta u + i V(x) u \).

Let's search for oscillating waves. For the Schr\"odinger equation, we have solutions \( u(x, t) = e^{i x \xi} e^{i t \omega(\xi)} \), which are oscillating waves so long as \( \omega \colon \mathbb{R}^{d} \to \mathbb{C} \) ends up having no imaginary part, this ends up being an oscillator. We want to find \( \xi \in \mathbb{R}^{d} \) and suitable \( \omega \) for this. So let's plug into our equation.
\[ \partial_{t} u = i \omega(\xi) \cdot u \quad i \Delta u = i \left( - | \xi|^{2} \right) \cdot u. \]
So if we have \( \omega( \xi) = - | \xi|^{2} \), we have an oscillating solution
\[ u(x, t) = e^{i x \xi} e^{-i t | \xi|^{2}}. \]

